\section{Results and Discussion}
\label{sec:results}

\subsection{RQ1: Does Trace-Complexity Bias Exist?}

The baseline results indicate that trace-complexity bias is substantial.
Table~\ref{tab:complexity-bias-baseline} shows that anomaly scores
produced by the baseline STV representation combined with Isolation
Forest are strongly correlated with trace structural complexity. Across
the three datasets, the Pearson correlation between anomaly score and
the number of services ranges from \(r=0.968\) to \(0.982\), while the
correlation with the number of spans ranges from \(r=0.854\) to
\(0.981\). These consistently high correlations indicate that
structurally larger traces receive systematically higher anomaly scores
even in the absence of injected anomalies. This demonstrates that trace
complexity introduces a substantial bias into the baseline anomaly
scoring process and provides the primary motivation for WR-STV.

\subsection{RQ2: Does Residual Normalization Reduce It?}

Residual normalization substantially reduces, but does not eliminate,
\begin{table}[t]
\centering
\caption{Complexity-bias correlation across methods, averaged over
num\_services and num\_spans (real, non-injected test traces). This is the
paper's central bias-reduction evidence.}
\label{tab:complexity-bias-all-methods}
\begin{tabular}{lcc}
\toprule
\textbf{Method} & \textbf{corr.\ num\_services} & \textbf{corr.\ num\_spans} \\
\midrule
Baseline STV + IF                 & 0.868 & 0.820 \\
Enhanced STV + IF                 & 0.835 & 0.779 \\
Global Residual STV (unweighted)  & 0.383 & 0.324 \\
\textbf{WR-STV (weighted)}        & \textbf{0.296} & \textbf{0.238} \\
Request-Aware WR-STV              & 0.274 & 0.218 \\
\bottomrule
\end{tabular}
\end{table}

trace-complexity bias. Table~\ref{tab:complexity-bias-all-methods} shows that
WR-STV considerably reduces the dependence of anomaly scores on trace
structural complexity. Averaged across datasets and seeds, the Pearson
correlation with the number of services decreases from \(r=0.868\)
under the baseline detector to \(r=0.296\) using WR-STV and further to
\(r=0.274\) under the request-aware variant. Similarly, the correlation
with the number of spans decreases from \(r=0.820\) to \(r=0.238\) and
\(r=0.218\), respectively.

The remaining positive correlation is expected because larger traces
naturally provide more opportunities for at least one extreme feature
residual to occur. Consequently, WR-STV should be interpreted as
substantially reducing, rather than completely removing,
trace-complexity bias while preserving sensitivity to path-visible
latency deviations.

\subsection{RQ3: Does WR-STV Improve Anomaly Scoring and How Does It Compare with Baselines?}
\label{sec:results-main}
\begin{table}[t]
\centering
\caption{Main detection results: baseline vs.\ WR-STV, largest-visible
injection ($f=30$, $m=3$), mean over datasets/seeds.}
\label{tab:main-results}
\begin{tabular}{lccc}
\toprule
& \textbf{ROC-AUC} & \textbf{AUPRC} & \textbf{F1} \\
\midrule
\multicolumn{4}{l}{\textit{Within-dataset}} \\
Baseline STV + IF & 0.605 & 0.574 & 0.583 \\
\textbf{WR-STV ($k$=1, $C$=5)} & \textbf{0.814} & \textbf{0.812} & \textbf{0.738} \\
\midrule
\multicolumn{4}{l}{\textit{Cross-dataset (6 directed pairs)}} \\
Baseline STV + IF & 0.622 & 0.573 & 0.602 \\
\textbf{WR-STV ($k$=1, $C$=5)} & \textbf{0.803} & \textbf{0.804} & \textbf{0.722} \\
\bottomrule
\end{tabular}
\end{table}

\begin{table}[t]
\centering
\caption{Paired statistical tests, WR-STV vs.\ baseline STV~+~IF.}
\label{tab:main-stats}
\begin{tabular}{lcccc}
\toprule
& \textbf{Gain} & \textbf{Cohen's $d_z$} & \textbf{paired-$t$ $p$} & \textbf{Wilcoxon $p$} \\
\midrule
\multicolumn{5}{l}{\textit{Within-dataset (n=15 seed$\times$dataset rows)}} \\
ROC-AUC & +0.209 & 2.78 & $3.68 \times 10^{-8}$ & $6.1 \times 10^{-5}$ \\
AUPRC   & +0.237 & --   & --                     & -- \\
F1      & +0.155 & --   & --                     & -- \\
\midrule
\multicolumn{5}{l}{\textit{Cross-dataset (n=6 directed pairs)}} \\
ROC-AUC & +0.182 & 1.92 & $0.0053$ & $0.03125$ \\
AUPRC   & +0.231 & --   & --       & $0.03125$ \\
F1      & +0.120 & --   & --       & $0.03125$ \\
\bottomrule
\end{tabular}
\end{table}


Table~\ref{tab:main-results} and Table~\ref{tab:main-stats} report the
primary anomaly scoring comparison. Compared with baseline STV +
Isolation Forest, WR-STV improves the within-dataset ROC-AUC from
\(0.605\) to \(0.814\). The cross-dataset ROC-AUC also improves from
\(0.622\) to \(0.803\). The within-dataset improvement is statistically
significant with a large paired effect size. The cross-dataset
comparison also shows a large paired effect, although its statistical
interpretation is limited by the small number of directed train-test
pairs.

These results indicate that the primary performance improvement arises
from feature-wise residual normalization rather than from replacing one
unsupervised detector with another. Standardizing each service-path
feature relative to its learned latency distribution produces a scoring
space that is substantially less influenced by trace structural
complexity, allowing latency deviations to be identified more
consistently across heterogeneous traces.
\begin{table}[t]
\centering
\caption{Detection performance vs.\ classical baselines, random-visible
injection ($f=30$), mean ROC-AUC.}
\label{tab:extra-baselines}
\begin{tabular}{lc}
\toprule
\textbf{Method} & \textbf{ROC-AUC} \\
\midrule
\textbf{WR-STV ($k$=1, $C$=5)}   & \textbf{0.811} \\
Max-Z Residual STV (unweighted) & 0.798 \\
Residual STV + LOF              & 0.785 \\
Residual STV + One-Class SVM    & 0.775 \\
Top-3 Residual STV (unweighted) & 0.771 \\
Raw STV + One-Class SVM         & 0.761 \\
Raw STV + LOF                   & 0.707 \\
Enhanced STV + IF               & 0.673 \\
Baseline STV + IF               & 0.604 \\
\bottomrule
\end{tabular}
\end{table}


Table~\ref{tab:extra-baselines} further compares WR-STV with classical
baselines operating on the same residual representation. Although WR-STV
achieves the highest overall ROC-AUC (\(0.811\)), the margin over the
unweighted Max-Z Residual STV (\(0.798\)) is relatively small.
\begin{table}[t]
\centering
\caption{Detection performance vs.\ classical baselines, random-visible
injection ($f=30$), mean ROC-AUC.}
\label{tab:extra-baselines}
\begin{tabular}{lc}
\toprule
\textbf{Method} & \textbf{ROC-AUC} \\
\midrule
\textbf{WR-STV ($k$=1, $C$=5)}   & \textbf{0.811} \\
Max-Z Residual STV (unweighted) & 0.798 \\
Residual STV + LOF              & 0.785 \\
Residual STV + One-Class SVM    & 0.775 \\
Top-3 Residual STV (unweighted) & 0.771 \\
Raw STV + One-Class SVM         & 0.761 \\
Raw STV + LOF                   & 0.707 \\
Enhanced STV + IF               & 0.673 \\
Baseline STV + IF               & 0.604 \\
\bottomrule
\end{tabular}
\end{table}

Table~\ref{tab:extra-baselines} shows that this difference is not
statistically significant for ROC-AUC (\(p=0.108\)) or F1
(\(p=0.735\)), while only Average Precision reaches statistical
significance. This suggests that most of WR-STV's performance gain
originates from residual normalization, whereas reliability weighting
provides a smaller, metric-dependent incremental improvement.

\subsection{RQ4: How Robust Is WR-STV to Injection Strength?}
\label{sec:results-sensitivity}
\begin{table}[t]
\centering
\caption{Sensitivity to injection strength $f$, largest-visible strategy,
WR-STV vs.\ baseline.}
\label{tab:sensitivity}
\begin{tabular}{lcccc}
\toprule
& \multicolumn{4}{c}{\textbf{Injection factor $f$}} \\
& \textbf{5$\times$} & \textbf{10$\times$} & \textbf{20$\times$} & \textbf{30$\times$} \\
\midrule
\multicolumn{5}{l}{\textit{WR-STV ROC-AUC / AUPRC / F1}} \\
& 0.646 / 0.625 / 0.603 & 0.724 / 0.711 / 0.660 & 0.784 / 0.780 / 0.715 & 0.814 / 0.812 / 0.738 \\
\midrule
Baseline STV+IF ROC-AUC & 0.581 & -- & -- & 0.605 \\
\bottomrule
\end{tabular}

\vspace{4pt}
\footnotesize
WR-STV improves monotonically across all three metrics from $5\times$ to
$30\times$; baseline STV+IF is comparatively insensitive to injection
strength (ROC-AUC 0.581 $\rightarrow$ 0.605 over the same range).
\todo{Fill the two baseline ROC-AUC blanks (10x, 20x) from
\texttt{sensitivity\_key\_df} if you want the full baseline row --
optional, the endpoints already make the monotonicity-vs-insensitivity
point.}
\end{table}

Table~\ref{tab:sensitivity} shows that WR-STV's ROC-AUC, AUPRC, and F1
increase monotonically as the injection factor increases from
\(5\times\) to \(30\times\), while baseline STV~+~IF is comparatively
insensitive to injection strength over the same range. This indicates
that WR-STV's score responds to the magnitude of controlled latency
deviations, which is desirable for an anomaly scoring method. However,
even the smallest evaluated factor, \(5\times\), represents a large
multiplicative latency deviation. These results therefore demonstrate
sensitivity trends under controlled perturbations rather than
performance under subtle production-scale degradations.

We additionally evaluate the random-visible injection strategy, which
selects injected spans uniformly at random among STV-visible spans rather
than always selecting the largest-duration visible span. This checks
whether WR-STV's advantage is an artifact of injections always landing
on spans most likely to affect a top-residual scoring rule. WR-STV's
ROC-AUC changes negligibly between strategies
(\(0.814 \rightarrow 0.811\)), as does the baseline's
(\(0.605 \rightarrow 0.604\)). This indicates that the observed
advantage is not explained by this specific form of injection
circularity. Nevertheless, both strategies remain restricted to
STV-visible service paths and model only multiplicative latency
anomalies, limitations discussed further in Section~\ref{sec:threats}.

\subsection{RQ5: Does WR-STV Generalize Across Datasets?}
\begin{table}[t]
\centering
\caption{WR-STV cross-dataset transfer, all six directed train$\rightarrow$test
pairs (largest-visible injection)}.
\label{tab:cross-pairs}
\begin{tabular}{lccc}
\toprule
\textbf{Train $\rightarrow$ Test} & \textbf{ROC-AUC} & \textbf{AUPRC} & \textbf{F1} \\
\midrule
Dataset 1 $\rightarrow$ Dataset 2 & 0.891 & 0.888 & 0.800 \\
Dataset 1 $\rightarrow$ Dataset 3 & 0.891 & 0.892 & 0.796 \\
Dataset 2 $\rightarrow$ Dataset 1 & 0.611 & 0.644 & 0.587 \\
Dataset 2 $\rightarrow$ Dataset 3 & \todo{fill} & \todo{fill} & \todo{fill} \\
Dataset 3 $\rightarrow$ Dataset 1 & \todo{fill} & \todo{fill} & \todo{fill} \\
Dataset 3 $\rightarrow$ Dataset 2 & 0.872 & \todo{fill} & \todo{fill} \\
\midrule
\textbf{Mean (6 pairs)} & \textbf{0.803} & \textbf{0.804} & \textbf{0.722} \\
\bottomrule
\end{tabular}
\end{table}


Table~\ref{tab:cross-pairs} reports all six directed train-test pairs.
Transfer is strong for five of the six pairs, with ROC-AUC values of at
least \(0.87\), but weak for Dataset~2~\(\rightarrow\)~Dataset~1
(ROC-AUC \(0.611\)). We checked whether this degradation is explained by
unseen service paths at test time and found that \(97.8\%\) of
test-trace STV-visible events were covered by the Dataset~2 training
vocabulary. This suggests that the weak transfer pair is not primarily a
coverage problem.
\begin{table}[t]
\centering
\caption{Latency distribution shift vs.\ cross-dataset transfer performance
(shared STV-visible features between train and test).}
\label{tab:distribution-shift}
\begin{tabular}{lccc}
\toprule
\textbf{Train $\rightarrow$ Test} & \textbf{Wtd.\ mean Wasserstein} & \textbf{$p_{90}$ std.\ shift} & \textbf{ROC-AUC} \\
\midrule
Dataset 2 $\rightarrow$ Dataset 1 & 0.973 & 2.737 & 0.611 \\
Dataset 1 $\rightarrow$ Dataset 3 & 0.872 & 1.515 & 0.891 \\
Dataset 1 $\rightarrow$ Dataset 2 & 0.772 & 1.707 & 0.891 \\
\bottomrule
\end{tabular}

\vspace{4pt}
\footnotesize
Spearman correlation ($p_{90}$ standardized shift vs.\ ROC-AUC, all 6
directed pairs): $r = -0.829$, $p = 0.0416$. \textbf{Report as exploratory}
-- $n=6$ pairs are drawn from only 3 underlying datasets and are not
independent samples; do not present this $p$-value as confirmatory.
\end{table}


Table~\ref{tab:distribution-shift} instead shows that this pair has the
largest latency distribution shift on shared features
(\(p_{90}\) standardized shift \(=2.74\), compared with
\(\leq 1.71\) for the stronger pairs shown). The exploratory correlation
across all six directed pairs (Spearman \(r=-0.83\), \(n=6\)) is more
consistent with distribution shift than with vocabulary coverage as an
explanation for transfer degradation. This suggests that successful
cross-dataset transfer may depend more strongly on preserving the
latency distributions of shared service paths than on maintaining
identical feature vocabularies. Since this conclusion is based on only
six directed transfer pairs, it should be regarded as exploratory rather
than definitive.

\subsection{RQ6: Can WR-STV Explain Anomalies?}
\label{sec:results-explainability}
\begin{table}[t]
\centering
\caption{WR-STV explanation overlap against known injected paths
(aggregate cross-dataset, largest-visible injection).}
\label{tab:explainability}
\begin{tabular}{lcc}
\toprule
\textbf{Top-$k$} & \textbf{Service overlap} & \textbf{Path overlap} \\
\midrule
Top-1  & 0.638 & 0.500 \\
Top-3  & 0.937 & 0.915 \\
Top-5  & 0.967 & 0.946 \\
Top-10 & 0.998 & 0.992 \\
\bottomrule
\end{tabular}
\end{table}

Table~\ref{tab:explainability} shows that WR-STV's top-3 contributing
service paths recover the actual injected service paths with
\(93.7\%\) service overlap and \(91.4\%\) path overlap, aggregated across
all cross-dataset pairs. This result is conditional on the injected
anomaly being among the largest residual signals in the trace. The same
scoring mechanism used for anomaly scoring is also used for localization,
so this result should be interpreted as evidence that WR-STV can point to
anomalies it scores highly, not as an independently validated
localization capability for anomalies of unknown type.

\begin{figure}[t]
\centering
% \includegraphics[width=0.45\textwidth]{figures/explanation_example.pdf}
\caption{Example WR-STV explanation for a single injected trace.}
\label{fig:explanation-example}
\end{figure}

\subsection{RQ7: Where Does WR-STV Fail?}
\begin{table}[t]
\centering
\caption{WR-STV failure taxonomy: cases where the injected trace did not
score above its original (279 unique failures after deduplication from
1{,}395 raw pairwise failure rows).}
\label{tab:failure-categories}
\begin{tabular}{lcc}
\toprule
\textbf{Category} & \textbf{Count} & \textbf{\%} \\
\midrule
Pre-existing high normal score       & 124 & 44.44 \\
Very short trace, limited evidence   & 105 & 37.63 \\
Score unchanged after injection      & 19  & 6.81 \\
Low residual after injection         & 12  & 4.30 \\
Uncategorized (distribution/ranking) & 10  & 3.58 \\
Low reliability weight               & 5   & 1.79 \\
High train-variance path             & 4   & 1.43 \\
\bottomrule
\end{tabular}
\end{table}

Table~\ref{tab:failure-categories} categorizes the 279 unique cases,
after deduplication from 1{,}395 raw pairwise rows, where WR-STV did not
rank an injected trace above its original. The two dominant categories,
pre-existing high normal score (\(44.4\%\)) and very short traces with
limited path evidence (\(37.6\%\), mean \(\approx 1.3\) spans), together
account for \(82\%\) of failures. Both categories have clear causal
interpretations: a trace that already scores in the upper tail of the
normal-score distribution has little room to be pushed higher by an
injected anomaly, and a one-or-two-span trace provides limited
STV-visible evidence for a top-\(k\) aggregation with \(k=1\).

These failure modes suggest that many WR-STV failures arise from limited
or unfavorable evidence available to the scorer rather than from random
instability of the scoring rule. In practice, traces with very few
observable service-path features may require confidence estimates,
abstention rules, or complementary detectors rather than being treated
as equally reliable scoring targets.

\subsection{RQ8: Is WR-STV Computationally Lightweight?}

Although computational efficiency is not the primary contribution of
this work, WR-STV also offers practical scoring advantages.
\begin{table}[t]
\centering
\caption{Runtime comparison, mean over datasets and seeds.}
\label{tab:runtime}
\begin{tabular}{lcc}
\toprule
\textbf{Method} & \textbf{Total (s)} & \textbf{Per-trace scoring (ms)} \\
\midrule
\textbf{WR-STV}    & \textbf{0.153} & \textbf{0.0238} \\
Enhanced STV + IF  & 0.537 & 0.3627 \\
Baseline STV + IF  & 0.551 & 0.3647 \\
\bottomrule
\end{tabular}

\vspace{4pt}
\footnotesize
WR-STV is $\approx$3.5$\times$ faster end-to-end and $\approx$15$\times$
faster per-trace than Isolation Forest-based baselines. Memory footprint is
negligible for all methods at this dataset scale and is not a
differentiator; report it in one sentence in text, not as a table.
\end{table}

Table~\ref{tab:runtime} shows that WR-STV requires approximately
\(0.022\)\,ms per trace, compared with approximately \(0.36\)\,ms for
Isolation Forest-based baselines, corresponding to roughly a
\(16\times\) reduction in inference time. This improvement is expected
because WR-STV performs lightweight statistical computations using
precomputed feature statistics and avoids ensemble-tree inference during
deployment. Consequently, WR-STV has favorable computational properties
for high-throughput scoring scenarios, although full online deployment
would require additional thresholding and operational validation.

\subsection{Summary of Findings}

Across the evaluated research questions, the results support the central
hypothesis of this work. Raw STV representations scored with Isolation
Forest are strongly influenced by trace structural complexity. WR-STV
substantially reduces this dependence through feature-wise residual
normalization and improves controlled latency-anomaly ranking in
within-dataset experiments, with promising but more limited evidence for
cross-dataset transfer. The method also provides service-path-level
interpretability, predictable failure modes, and lightweight scoring
cost. Overall, the results support WR-STV as a bias-aware anomaly
scoring framework for controlled STV-visible latency deviations rather
than a general-purpose solution to all microservice anomaly types.
